\subsection{Navigation}

\subsubsection{Stack overview.}
FRIDA's navigation is built on Nav2~\cite{macenski2020marathon2}.
SLAM Toolbox builds 2D maps of the arena from
the combined scan of the two RPLIDAR~C1 sensors, using graph-based SLAM with
loop closure, and is then run in localization mode on the saved map. Wheel and
inertial odometry are fused by an EKF on the base controller
and provide the motion estimate for localization.
\subsubsection{Planning and Control.}
Global paths are computed with Nav2's NavFn planner using A* search, while a
Model Predictive Path Integral (MPPI) controller~\cite{williams2016mppi}
handles local trajectory tracking. The controller uses a holonomic motion
model, allowing FRIDA to move laterally when approaching furniture or
navigating narrow passages. The local and global costmaps combine the lidar
scan with a spatio-temporal voxel layer built from the RGB-D point cloud,
which captures obstacles outside the lidar plane and lets stale voxels decay
over time, so that phantom obstacles are cleared even while the robot is
stationary.
\subsubsection{Behavior Tree and Supervision.}
Nav2's behavior tree coordinates path planning, goal execution and recovery
behaviors. A custom supervisor monitors the sensor streams and transforms
required by SLAM and Nav2. If any of them becomes unavailable, it pauses both
components and resumes them once the required data is restored. The supervisor
also provides the task manager with a reduced set of navigation services.
